\begin{frame}[shrink]
	\frametitle{Treinar através de uma função que não tem gradiente}
	\begin{itemize}
		\item $\text{round}(\cdot)$ e $\text{clip}(\cdot)$ são \textbf{funções em degrau}: a derivada é zero em quase todo ponto, e indefinida nos degraus.
		\item Retropropagação padrão não tem o que propagar através delas --- o gradiente morre.
	\end{itemize}
	\vspace{1pt}
	\begin{block}{A solução: \textit{Straight-Through Estimator} (STE)}
		\begin{itemize}
			\item \textbf{Forward}: usa os pesos quantizados $\widetilde{W}$ normalmente.
			\item \textbf{Backward}: finge que a quantização foi a função identidade --- deixa o gradiente passar direto para uma cópia dos pesos em precisão plena, mantida ``nos bastidores'' e só quantizada de novo a cada passada.
		\end{itemize}
	\end{block}
	\vspace{0pt}
	\begin{alertblock}{Ponte com a Parte II}
		É a \textbf{mesma} ideia por trás do \textit{gradiente substituto} em redes de pulso \cite{eshraghian2023training}: forward com a função em degrau real, backward com uma aproximação suave, diferenciável. Duas subáreas, independentemente, chegaram à mesma solução para o mesmo problema estrutural.
	\end{alertblock}
\end{frame}
